% source: https://github.com/pfradin/luadraw/discussions/228#discussioncomment-16130420 (pfradin, block 1)
\documentclass{standalone}
\usepackage{luadraw}
\usepackage{fourier}
\begin{document}

\begin{luadraw}{name=cosine_graph}
    local pi, i = math.pi, cpx.I
    local g = graph:new{
        window={-2.5, 7, -7, 4.5}, size={12, 9, 8/pi},
    }
    g:Linewidth(8) -- for axes thickness
    g:Daxes(
        {0, pi/2, 2}, 
        {
            limits = {{-1.85, 7}, {-5.75, 4.3}},
            gradlimits = {{-1.85, 6.5}, {-5.75, 3.5}},
            grid = true, gridcolor="black", 
            nbsubdiv={ 3, 1 }, subgridcolor="black",
            arrows = "-latex",
            xyticks = { 0, 0 },
            labeltext = {"\\pi",""},
            labelden = { 2, 1 },
            legend = {"$x$", "$y$"},
            legendpos = { 1, 1 },
            legendsep = { 0, 0 },
            xynode_options = "fill=white,inner sep=1.75pt",
            labelshift = { 0 , -0.1875 }, -- Now labelshift takes precedence over defaultlabelshift
            xylabelsep = { 0, 0 },
            originloc=i, originnum={0,1}, -- originloc is now the origin point for graduations (not the intersection of the axes)
            originpos = {"none", ""}, -- the origin graduation on y axis is 1 (not 0)
        }
    )

    local f = function(x) return -3*math.cos(2*x)-1 end
    g:Dcartesian(f, {
        x = {-pi/2 - 0.2, 2*pi + 0.15},
        nbdots = 100, draw_options = "line width=1.2pt"
    })
    g:Dlabel("$O$", 0, {pos="SW"})
    g:Dlabel("Graph of $f$", Z(2.35, -6.5), {})
    g:Show()
\end{luadraw}
\end{document}
